% !TeX root = ../main.tex

\chapter{Datos y Metodología}

\section{El Catálogo de Hunt y Reffert}

    Este trabajo utiliza como base fundamental el catálogo de cúmulos estelares 
    desarrollado por Hunt y Reffert, el cual es el resultado de un esfuerzo 
    sistemático por actualizar y limpiar el censo de agrupaciones estelares en 
    la Vía Láctea utilizando aprendizaje automático no supervisado. El catálogo 
    final es el producto de una metodología desarrollada a lo largo de tres 
    estudios consecutivos, evolucionando desde la selección de algoritmos de 
    agrupamiento espacial, pasando por la implementación a nivel de todo el 
    cielo, hasta la limpieza dinámica del censo basada en masas y radios de 
    Jacobi.

    \subsection{Descripción de los datos}

        El censo se construye utilizando los datos astrométricos y fotométricos 
        del tercer conjunto de datos de la misión Gaia (Gaia DR3). A diferencia 
        de catálogos anteriores que operaban con el catálogo Gaia DR2 y cortes 
        de magnitud muy restrictivos ($G \le 18$), el catálogo definitivo de 
        Hunt y Reffert procesa una base de aproximadamente 729.7 millones de 
        fuentes estelares.

        Los autores utilizan el espacio de cinco dimensiones (5D) proporcionado 
        por la astrometría de Gaia para la búsqueda de cúmulos:

        \begin{itemize}
            \item \textbf{Posiciones espaciales:} Ascensión recta ($\alpha$) y 
            declinación ($\delta$).

            \item \textbf{Cinemática:} Movimientos propios en ambas coordenadas 
            espaciales ($\mu_{\alpha*}$, $\mu_\delta$).

            \item \textbf{Distancia:} Paralaje ($\varpi$).
        \end{itemize}

        Para evitar distorsiones esféricas e inhomogeneidades en la densidad de 
        los datos, el cielo se divide utilizando la teselación HEALPix en 
        diferentes niveles según la profundidad de la búsqueda, y los datos de 
        las cinco dimensiones son re-escalados y centrados antes de aplicar los 
        algoritmos de agrupamiento.

    \subsection{Criterios de Selección y Preprocesamiento}
        La metodología de selección de miembros y validación de cúmulos 
        realizada por los autores se destaca por aplicar rigurosos filtros de 
        calidad, tanto a nivel de estrellas individuales como a nivel de la 
        agrupación candidata:

        \begin{enumerate}
            \item \textbf{Corte de calidad astrométrica:} En lugar de aplicar un 
            corte estricto de magnitud, se utiliza una red neuronal indicadora 
            de calidad astrométrica (específicamente la bandera de calidad v1 de 
            Rybizki et al. con un valor $> 0.5$). Esto permite recuperar 
            estrellas miembro más débiles de hasta magnitudes $G \sim 20$ que 
            aún poseen astrometría confiable, aumentando drásticamente la 
            relación señal-ruido de los cúmulos distantes.

            \item \textbf{Algoritmo de agrupamiento:} El catálogo emplea el 
            algoritmo HDBSCAN, seleccionado por su sensibilidad superior para 
            detectar cúmulos en conjuntos de datos con densidades variables.

            \item \textbf{Validación fotométrica:} Las agrupaciones detectadas 
            son filtradas mediante Redes Neuronales Convolucionales Bayesianas 
            entrenadas con isocronas PARSEC sintéticas, para confirmar que el 
            Diagrama Color-Magnitud (CMD) de la agrupación es estadísticamente 
            compatible con una población estelar simple y co-evolutiva.

            \item \textbf{Criterio de ligadura gravitacional:} Para el catálogo 
            final y ``limpio'', los autores filtran grupos de movimiento no 
            ligados (asociaciones) calculando la masa fotométrica total de cada 
            agrupación y derivando su Superficie de Roche o Radio de Jacobi 
            ($r_\text{J}$). Solo se retienen en la muestra de alta calidad 
            aquellos cúmulos con una probabilidad mayor al 50\% de poseer un 
            radio de Jacobi válido ($P(r_\text{J}) \ge 0.5$) y una masa mínima 
            ligada de $40\ \text{M}_\odot$. Esto redujo el censo drásticamente, 
            garantizando que los objetos de estudio son auténticos cúmulos 
            abiertos unidos gravitacionalmente.
        \end{enumerate}

    \subsection{Parámetros disponibles}
        
        El catálogo provee una parametrización exhaustiva que resulta ideal para 
        el análisis dinámico y estructural, proporcionando las siguientes 
        variables fundamentales para cada cúmulo:

        \begin{itemize}
            \item \textbf{Astrométricos y cinemáticos:} Posiciones medias, 
            paralajes medios, y movimientos propios medios derivados de los 
            miembros de alta probabilidad.

            \item \textbf{Parámetros Estructurales:} Radio que contiene el 50\% 
            de los miembros ($r_{50}$), radios del núcleo ($r_\text{c}$), radios 
            de marea aparentes ($r_\text{t}$), y radios de Jacobi 
            ($r_\text{J}$).

            \item \textbf{Parámetros Físicos y Fotométricos:} Edad ($\log t$), 
            extinción interestelar ($A_\text{V}$), extinción diferencial 
            ($\Delta A_\text{V}$) y módulo de distancia ($m-M$), todos inferidos 
            utilizando técnicas de Inferencia Variacional.
        \end{itemize}

\section{Tratamiento de Datos}

        \subsection{Criterios de Selección y Limpieza de la Muestra}

            El conjunto de datos inicial para este trabajo se obtuvo a partir 
            del catálogo de Hunt y Reffert, el cual consta originalmente de 7167 
            cúmulos estelares y 1291929 estrellas miembro. Dado que los archivos 
            originales de este catálogo se encuentran en formato FITS, se empleó 
            la librería Astropy para la lectura y extracción inicial de las 
            tablas.

            El catálogo original clasifica las estructuras detectadas en 
            diversas categorías mediante la columna \texttt{Type}, tales como 
            Open Clusters (OC), Globular clusters (GC), Moving groups (MG), Too 
            distant y Rejected. Mediante la aplicación de una máscara lógica 
            sobre esta columna, la muestra de trabajo se restringió 
            exclusivamente a los cúmulos abiertos (OC) y los cúmulos globulares 
            (GC). Una vez aislados los tipos de cúmulos de interés, se 
            identificaron sus respectivas estrellas miembro realizando un cruce 
            de datos mediante la columna \texttt{Name}, presente tanto en la 
            tabla de cúmulos como en la de miembros.

            Posteriormente, se llevó a cabo un proceso de limpieza del catálogo 
            de miembros para eliminar entradas duplicadas, utilizando la 
            librería Pandas y tomando como identificador único la columna 
            \texttt{GaiaDR3}. Para optimizar la manipulación de los datos en las 
            fases posteriores y aprovechar las herramientas de análisis de 
            Pandas, las tablas resultantes se exportaron a formato CSV. Durante 
            esta conversión, fue necesario realizar un preprocesamiento de las 
            cadenas de caracteres (strings) en las columnas \texttt{Name} y 
            \texttt{Type}, las cuales se importaban inicialmente como objetos de 
            bytes (\texttt{b'...'}); estos caracteres adicionales fueron 
            suprimidos para garantizar la correcta interpretación de las 
            etiquetas.

            Con el objetivo de realizar los cálculos estructurales adecuados, se 
            incorporaron a la base de datos las distancias fotométricas de 
            Bailer-Jones, disponibles en el archivo de Gaia. Estas distancias, 
            junto con sus respectivas incertezas, fueron extraídas mediante 
            consultas ADQL utilizando el identificador de Gaia. La integración 
            de esta información espacial geométrica al catálogo principal de 
            miembros se realizó mediante una unión en Pandas usando la columna 
            \texttt{GaiaDR3} en común, almacenando el resultado final nuevamente 
            en formato CSV.

            Para asegurar que el análisis de la dinámica interna de los sistemas 
            contará con suficiente significancia estadística, se impuso un 
            criterio de selección estricto, requiriendo un límite inferior de 
            mil miembros por cúmulo. Este criterio inicial redujo la muestra a 
            114 cúmulos masivos. Finalmente, se realizó una inspección 
            morfológica de estos sistemas, descartando 12 cúmulos que no 
            presentaban una simetría esférica suficiente, condición fundamental 
            para el posterior ajuste de los modelos teóricos de densidad. De 
            este modo, se consolidó una muestra final de 102 cúmulos estelares 
            para el análisis.

        \subsection{Sistemas de Referencia y Transformaciones}

            Durante el desarrollo del flujo de trabajo, se conservaron en su 
            mayoría las unidades físicas estándar empleadas en el catálogo 
            original: grados ($\text{deg}$) para coordenadas espaciales, 
            milisegundos de arco ($\text{mas}$) para paralajes, milisegundos de 
            arco por año ($\text{mas/yr}$) para movimientos propios, magnitudes 
            ($\text{mag}$) para fotometría, kilómetros por segundo 
            ($\text{km/s}$) para velocidades cinemáticas, masas solares 
            ($\text{M}_\odot$) y pársecs ($\text{pc}$) para distancias físicas. 
            Como única excepción, y con el fin de garantizar la compatibilidad 
            operativa con las funciones trigonométricas de la librería NumPy, 
            las coordenadas angulares se convirtieron sistemáticamente de grados 
            ($\text{deg}$) a radianes ($\text{rad}$).

            El sistema de referencia base adoptado corresponde al Sistema 
            Internacional de Referencia Celeste (ICRS) en la época J2016.0, 
            provisto por el catálogo. Sin embargo, para el análisis de cada 
            estructura individual, se definió un sistema de coordenadas relativo 
            centrado en cada cúmulo. El origen temporal de este sistema local se 
            hizo coincidir con el punto de mayor densidad del cúmulo, valor 
            reportado previamente en el catálogo de Hunt y Reffert. Esta 
            traslación del origen de coordenadas es de vital importancia, ya que 
            facilita directamente el cálculo de las distancias radiales de cada 
            estrella respecto al centro, paso indispensable para la construcción 
            de los perfiles radiales de densidad.